\documentclass[french]{book}
\usepackage{teach}
\usepackage{verbatim}
\usepackage{amsmath,amsfonts,amssymb}
\usepackage{pgfplots}
\usepackage{mwe}
\RequirePackage{graphicx}
\RequirePackage{ifpdf}
 \ifpdf
   \DeclareGraphicsRule{*}{mps}{*}{}
 \fi




\RequirePackage{fancyhdr}
\RequirePackage{titlesec}

  \usepackage{fancyhdr}

  \usepackage{amsmath}
  \usepackage{amssymb}
  \usepackage{amsfonts}
  \usepackage{amsthm}
  \usepackage{mathrsfs}

  \usepackage{array}
  \usepackage{multicol}
  \setlength{\multicolsep}{3pt}

  \usepackage{graphicx}
  \graphicspath{{Figures/}}
  \usepackage{pstricks,pst-plot,pst-text,pst-tree,pst-eucl}
  \usepackage[all]{xy}

  \usepackage{fancybox}
  \usepackage{color}

%  \usepackage[frenchb]{babel}
  \usepackage{hyperref}

\usepackage{subfig}
\pgfplotsset{compat=newest}
\usepackage[all]{xy}
%\usepackage{geometry}
%\geometry{top=1cm, bottom=1cm, left=2cm, right=2cm}
\usepackage[utf8]{inputenc}
\usepackage[french]{babel}
\redefinecolor{thm}{title}{bgcolor}{alizarin}
\redefinecolor{prop}{title}{bgcolor}{meatbrown}
\redefinecolor{defn}{title}{bgcolor}{olivine}
\definecolor{alizarin}{rgb}{0.82, 0.1, 0.26}
\definecolor{meatbrown}{rgb}{0.9, 0.72, 0.23}
\definecolor{olivine}{rgb}{0.6, 0.73, 0.45}
\usepackage{mathrsfs}
%%
\newcommand{\R}{\mathbb {R}}
%\newcommand{\C}{\mathds {C}}
\newcommand{\N}{\mathbb {N}}
\newcommand{\Z}{\mathbb {Z}}
\newcommand{\Q}{\mathbb {Q}}
\newcommand{\D}{\mathbb {D}}
%%
\newcommand{\se}{\geq}
\newcommand{\ie}{\leq}
\newcommand{\tm}{\times}
%\newcommand{\ell}{l}
%\newcommand{\ell'}{l'}
%%
\newcommand{\limn}{\lim_{n\rightarrow +\infty}}
\newcommand{\limo}[1][x]{\ds\lim_{#1\rightarrow 0}}
\newcommand{\limite}[2][x]{\ds\lim_{#1\rightarrow #2}}
\newcommand{\limpinf}[1][x]{\ds\lim_{#1\rightarrow +\infty}}
\newcommand{\limminf}[1][x]{\ds\lim_{#1\rightarrow -\infty}}
\newcommand{\limiteg}[2][x]{\ds\lim_{#1\xrightarrow{<}#2}}
\newcommand{\limited}[2][x]{\ds\lim_{#1\xrightarrow{>}#2}}
%%
\newcommand{\vect}[1]{\overrightarrow{#1}}
\newcommand{\oij}{(\textit{O},{\overrightarrow{i}},{\overrightarrow{j}})}
%
\newcommand{\cp}[2]{%
        \begin{pmatrix}
        #1\\
        #2
        \end{pmatrix}%
        }
%
\newcommand{\calb}{\mathscr{B}}
\newcommand{\calc}{\mathscr{C}}
\newcommand{\cald}{\mathscr{D}}
\newcommand{\cale}{\mathscr{E}}
\newcommand{\calf}{\mathscr{F}}
\newcommand{\calp}{\mathscr{P}}
\newcommand{\cals}{\mathscr{S}}
\newcommand{\calt}{\mathscr{T}}
%
\newcommand{\suite}[1][u]{\left( #1_{n} \right)_{n\in\N}}
\newcommand{\suitep}[1][u]{\left( #1_{n} \right)_{n\in\N^{*}}}
\usetikzlibrary{arrows}


\begin{document}
\tableofcontents
\chapter{Généralités sur les fonctions}
\section{Définition d'une fonction}
\begin{defn}
Une fonction $f$ est un objet mathématique qui est définie sur un ensemble de nombre de départ (appelé ensemble de définition) $D_f$ à valeurs dans un ensemble de nombre d'arrivé $F$. Elle associe à chaque élément de $E$ une valeur (appelé image) qui appartient à $F$. On note parfois la fonction $f$  comme ceci:  
$$f : E \longrightarrow F $$
$$ x \mapsto f(x)$$
\end{defn}
\section{Opérations sur les fonctions}
\subsection{Égalité de deux fonctions}
\begin{defn}
Soient $u$ et $v$ deux fonctions. On dit que $u$ et $v$ sont égales et on note $u=v$ si :
\begin{itemize}
\item $u$ et $v$ ont le même ensemble de définition $D$.
\item Pour tout $x\in D$, $u(x)=v(x)$.
\end{itemize}
\end{defn}

\begin{exemple}
Les fonctions $u$ et $v$ sont-elles égales ?
\begin{enumerate}
\item $u$ et $v$ sont définies par $u(x)=3-\dfrac{2}{x+1}$ et $v(x)=\dfrac{3x+1}{x+1}$
\item $u$ et $v$ sont définies par $u(x)=x$ et $v(x)=\dfrac{x^{2}}{x}$
\end{enumerate}
\end{exemple}


\begin{enumerate}
\item $u$ et $v$ ont le même ensemble de définition : $\R \backslash \{ -1\}$

Pour tout $x\in \R \backslash \{ -1\}$, \[u(x)=3-\dfrac{2}{x+1}=\dfrac{3(x+1)-2}{x+1}=\dfrac{3x+1}{x+1}=v(x)\]

donc $u=v$.
\item $u$ est définie sur $\R$ et $v$ est définie sur $\R^{*}$ donc $u\neq v$.
\end{enumerate}

%
\subsection{Opérations sur les fonctions}
\begin{defn}
Soient $u$ et $v$ deux fonctions définies sur $D$ et $\lambda$ un réel.
\begin{itemize}
\item On définit les fonctions $u+v$, $uv$, $\lambda u$, $u+\lambda$ de la façon suivante :
\[\begin{array}{l}
(u+v)(x)=u(x)+v(x) \\
(\lambda u)(x)=\lambda \times u(x)
\end{array}
\hspace{3em}
\begin{array}{l}
(uv)(x)=u(x)\times v(x) \\
(u+\lambda)(x)=u(x)+\lambda
\end{array}\]
\item Si, pour tout $x \in D$, $v(x)\neq 0$ alors on peut définir la fonction $\dfrac{u}{v}$ par :
\[\left(\dfrac{u}{v}\right)(x)=\dfrac{u(x)}{v(x)}\]

\end{itemize}
\end{defn}

\begin{exemple}
Soit $u$ et $v$ les fonctions définies sur $\R$ par $u(x)=x^{2}$ et $v(x)=x+3$. Déterminer $u+v$, $uv$, $2u$, $u+2$ et $\dfrac{u}{v}$.
\end{exemple}


\begin{itemize}
\item Pour tout réel $x$, $(u+v)(x)=x^{2}+x+3$ ; $(uv)(x)=x^{2}(x+3)=x^{3}+3x^{2}$ ; $(2u)(x)=2x^{2}$ et $(u+2)(x)=x^{2}+2$
\item Pour tout réel $x \neq -3$, $\left(\dfrac{u}{v}\right)(x)=\dfrac{x^{2}}{x+3}$
\end{itemize}

%
\subsection{Composition de fonctions}
\begin{defn}
Soit $u$ une fonction définie sur $D_{u}$ et $v$ une fonction définie sur $D_{v}$ et telle que pour tout $x\in D_{v}$, $v(x)\in D_{u}$.\\

Autrement dit on a:

\begin{minipage}[c]{0.5\linewidth}
\begin{center}
$$v : D_v \longrightarrow F_v (\subset D_u) $$
$$ x \mapsto v(x)$$
\end{center}
\end{minipage}
\hfill
\begin{minipage}[c]{0.5\linewidth}
\begin{center}
$$u : D_u \longrightarrow \mathbb{R} $$
$$ x \mapsto u(x)$$
\end{center}
\end{minipage}\\


On appelle fonction composée de $v$ par $u$ la fonction notée $u\circ v$ et définie sur $D_{v}$ par :\\

\hspace*{4.5cm}Pour tout $x\in D_{v}$, $u\circ v(x)=u(v(x))$


%\vspace{3ex}

\[\shorthandoff{;:!?} \xymatrix @!0 @R=0.8cm @C=2cm{
D_v \ar@{->}[r] & D_u \ar@{->}[r] & \R \\
x \ar@{|->}[r]^-v \ar@{|->}'[d]+D '[rrd]+D^-{u\circ v} [rr]& v(x) \ar@{|->}[r]^-u & u(v(x)) \\
& &
 }\]
\end{defn}
\vspace{1ex}

\begin{rem} Il faut faire attention à l'ordre des fonctions. $u\circ v$ et $v \circ u$ sont en général des fonctions différentes.
Il se peut qu'elles aient des ensembles de définition différents voire que l'une existe mais pas l'autre.
\end{rem}

\begin{exemple}
Soit $f$ la fonction définie sur $[0;+\infty[$ par $f(x)=\sqrt{x}-1$ et $g$ la fonction définie sur $\R$ par $g(x)=x^{2}+3$. Définir $g \circ f$ et $f \circ g$. Sont-elles égales ?
\end{exemple}


\begin{itemize}
\item $g \circ f$ est définie sur $[0;+\infty[$ par $g \circ f(x)=g(f(x))=(\sqrt{x}-1)^{2}+3=x-2\sqrt{x}+4$
\item $f \circ g$ est définie sur $\R$ par $f \circ g(x)=f(g(x))=\sqrt{x^{2}+3}-1$
\item $g \circ f$ et $f \circ g$ ne sont pas égales car elles n'ont pas le même ensemble de définition.
On peut ausi remarquer que $g \circ f(0) \neq f \circ g(0)$
\end{itemize}

%%%%%%%%%%%%%%%%%%%%%%%%%%%%%%%%%%%%%%%%%%%%%%%%%%%%%%%%%%%%%%%%%%%%%%%%%%%%%%%%%%%%%
\section{Sens de variation}
\subsection{Sens de variation de la fonction $u+\lambda$}
\begin{prop}
Soit $u$ une fonction défine sur un intervalle $I$ et $\lambda$ un réel.

Si $u$ est monotone sur $I$ alors $u$ et $u+\lambda$ ont même sens de variation sur $I$.
\end{prop}


\begin{demo}
Cas où $u$ est croissante

Soient $a$ et $b$ deux réels de $I$.
\begin{align*}
 a \ie b &\Longrightarrow u(a) \ie u(b) \text{ car $u$ est croissante sur $I$.}\\
  &\Longleftrightarrow u(a)+\lambda \ie u(b)+\lambda
\end{align*}

La fonction $u+\lambda$ est croissante sur $I$.

%\vspace{-\baselineskip}
\end{demo}
%
\subsection{Sens de variation de la fonction $\lambda u$}
\begin{prop}
Soit $u$ une fonction défine et monotone sur un intervalle $I$ et $\lambda$ un réel.
\begin{itemize}
\item Si $\lambda >0$ alors les fonctions $u$ et $\lambda u$ ont même sens de variation sur $I$.
\item Si $\lambda <0$ alors les fonctions $u$ et $\lambda u$ ont des sens de variation contraires sur $I$.
\end{itemize}
\end{prop}

%\ifthenelse{\value{version}>0}{\newpage}{}
\begin{demo}
Cas où $u$ est croissante

Soient $a$ et $b$ deux réels de $I$.
Si $a \ie b$ alors $u(a) \ie u(b)$ car $u$ est croissante sur $I$.
\begin{itemize}
\item Si $\lambda >0$ alors $\lambda u(a) \ie \lambda u(b)$ donc $\lambda u$ est croissante sur $I$.
\item Si $\lambda <0$ alors $\lambda u(a) \se \lambda u(b)$ donc $\lambda u$ est décroissante sur $I$.
\end{itemize}
\end{demo}
%
\subsection{Sens de variation de la fonction $u\circ v$}
\begin{prop}
Soit $u$ une fonction définie et monotone sur un intervalle $J$. Soit $v$ une fonction définie et monotone sur un intervalle $I$ et telle que pour tout $x\in I$, $v(x)\in J$.
\begin{itemize}
\item Si $u$ et $v$ ont même sens de variation alors $u\circ v$ est croissante sur $I$.
\item Si $u$ et $v$ ont des sens de variation contraires alors $u\circ v$ est décroissante sur $I$.
\end{itemize}
\end{prop}

\begin{demo}
\underline{Cas où $v$ est croissante sur $I$.}\\

Soient $a$ et $b$ deux réels de $I$.
Si $a \ie b$ alors $v(a) \in J$, $v(b) \in J$ et $v(a) \ie v(b)$ car $v$ est croissante sur $I$.
\begin{itemize}
\item Si $u$ est croissante sur $J$ alors $u(v(a)) \ie u(v(b))$ donc $u\circ v(a) \ie u\circ v(b)$ donc $u\circ v$ est croissante sur $I$.
\item Si $u$ est décroissante sur $J$ alors $u(v(a)) \se u(v(b))$ donc $u\circ v(a) \se u\circ v(b)$ donc $u\circ v$ est décroissante sur $I$.
\end{itemize}

\underline{Cas où $v$ est décroissante sur $I$.} \\

Soient $a$ et $b$ deux réels de $I$.
Si $a \ie b$ alors $v(a) \in J$, $v(b) \in J$ et $v(a) \se v(b)$ car $v$ est décroissante sur $I$.
\begin{itemize}
\item Si $u$ est croissante sur $J$ alors $u(v(a)) \se u(v(b))$ donc $u\circ v(a) \se u\circ v(b)$ donc $u\circ v$ est décroissante sur $I$.
\item Si $u$ est décroissante sur $J$ alors $u(v(a)) \ie u(v(b))$ donc $u\circ v(a) \ie u\circ v(b)$ donc $u\circ v$ est croissante sur $I$.
\end{itemize}

\end{demo}
%%%%%%%%%%%%%%%%%%%%%%%%%%%%%%%%%%%%%%%%%%%%%%%%%%%%%%%%%%%%%%%%%%%%%%%%%%%%%%%%%%%%%
\newpage
\section{Représentations graphiques}

\subsection{Définition d'une courbe}
\begin{defn}
La courbe représentative de la fonction $f$, notée $C_f$,  est l’ensemble des points de coordonnées $(x ; f(x))$, avec $x$ qui prend \emph{(une par une)} toutes les valeurs de l'ensemble $D_f$. \\
 
Autrement dit, $$ C_f= \{ (x ; f(x))\; /  \; x\in D_f \} $$
\end{defn}

\subsection{Représentation graphique d'une fonction $x \mapsto u(x+a)+b$}
\begin{prop}
Soit $u$ une fonction et $v$ la fonction définie par $v(x)=u(x+a)+b$.

Dans un repère $\oij$, on appelle $\calc_{u}$ et $\calc_{v}$ les courbes représentatives des fonctions $u$ et $v$.

$\calc_{v}$ est l'image de $\calc_{u}$ par la translation de vecteur $-a\vect{i}+b\vect{j}$, autrement dit le vecteur de coordonnées $\cp{-a}{b}$.
\end{prop}
\begin{center}
\includegraphics{Figures/FiguresFonctionsCours.1}
\end{center}

%\ifthenelse{\value{version}>0}{\newpage}{}
\begin{demo}
Soient $M(x;y)$ et $M'(x-a;y+b)$.

\[M'\in \calc_{v} \Leftrightarrow y+b=v(x-a) \Leftrightarrow y+b=u(x-a+a)+b \Leftrightarrow y=u(x) \Leftrightarrow M \in \calc_{v}\]
\end{demo}


\subsection{Représentation graphique d'une fonction $\lambda u$}
\begin{prop}
Soit $u$ une fonction et $v$ la fonction $\lambda u$
Dans un repère $\oij$, on appelle $\calc_{u}$ et $\calc_{v}$ les courbes représentatives des fonctions $u$ et $v$.
Si $M$ est le point de $\calc_{u}$ d'abscisse $x$ alors on obtient le point d'abscisse $x$ de $\calc_{v}$ en multipliant l'ordonnée de $M$ par $\lambda$.
\end{prop}

\begin{minipage}[c]{0.3\linewidth}
\begin{center}
\includegraphics[scale=0.9]{Figures/FiguresFonctionsCours.2}
\end{center}
\end{minipage}
\hfill
\begin{minipage}[c]{0.3\linewidth}
\begin{center}
\includegraphics[scale=0.9]{Figures/FiguresFonctionsCours.3}
\end{center}
\end{minipage}
\hfill
\begin{minipage}[c]{0.3\linewidth}
\begin{center}
\includegraphics[scale=0.9]{Figures/FiguresFonctionsCours.4}
\end{center}
\end{minipage}
\end{document}
